\section{Topic-Level Results}
\label{sec:topic-level-results}

Having established that community-level aggregation reveals paradigm-specific effects (Section~\ref{sec:community-level-results}), we now present the finer-grained topic-level analysis for completeness. These results are also available for five annotation paradigms: Human, Fine-tuned RoBERTa, GPT-4o-mini (baseline), Fine-tuned GPT-4o-mini, and Llama-3.3-70B-versatile.

\subsection{Mediation Analysis Results}
\label{sec:mediation-results}

We evaluate \textbf{Donald Trump sentiment as treatment} with Politics sentiment as mediator ($N = 233$ articles); the relatively small $N$ reflects the sparsity of Trump-tagged articles in the corpus.

Table~\ref{tab:te-trump} reports Total Effect estimates for both percentile contrasts across the five annotation paradigms. The treatment window is identical across all four datasets ($Q25 = -0.7$, $Q75 = +0.1$; $P10 = -1.1$, $P90 = +0.3$), so differences across paradigms arise entirely from differences in ideology labels $Y$.

\begin{table}[ht]
  \centering
  \sisetup{table-format=+1.3}
  \begin{tabular}{llS[table-format=+1.3]S[table-format=+1.3]S[table-format=+1.3]S[table-format=+1.3]S[table-format=+1.3]}
    \toprule
    \textbf{Paradigm} & \textbf{Contrast} &
      {\textbf{TE}} &
      {\textbf{CI lower}} & {\textbf{CI upper}} &
      {\textbf{NDE}} & {\textbf{NIE}} \\
    \midrule
    Human
      & $Q25 \to Q75$  & -0.199 & -0.426 & +0.027 & -0.077 & -0.123 \\
    Human
      & $P10 \to P90$  & -0.344 & -0.735 & +0.047 & -0.120 & -0.223 \\
    \addlinespace
    Fine-tuned RoBERTa
      & $Q25 \to Q75$  & -0.103 & -0.323 & +0.117 & -0.034 & -0.069 \\
    Fine-tuned RoBERTa
      & $P10 \to P90$  & -0.177 & -0.557 & +0.202 & -0.047 & -0.130 \\
    \addlinespace
    GPT-4o-mini (baseline)
      & $Q25 \to Q75$  & -0.078 & -0.298 & +0.143 & +0.003 & -0.080 \\
    GPT-4o-mini (baseline)
      & $P10 \to P90$  & -0.134 & -0.515 & +0.247 & -0.003 & -0.131 \\
    \addlinespace
    Fine-tuned GPT-4o-mini
      & $Q25 \to Q75$  & {$\mathbf{-0.249^{*}}$} & {$\mathbf{-0.465}$} & {$\mathbf{-0.033}$} & -0.101 & -0.148 \\
    Fine-tuned GPT-4o-mini
      & $P10 \to P90$  & {$\mathbf{-0.430^{*}}$} & {$\mathbf{-0.802}$} & {$\mathbf{-0.058}$} & -0.168 & -0.262 \\
    \addlinespace
    Llama-3.3-70B
      & $Q25 \to Q75$  & -0.091 & -0.338 & +0.157 & -0.021 & -0.070 \\
    Llama-3.3-70B
      & $P10 \to P90$  & -0.156 & -0.583 & +0.270 & -0.041 & -0.116 \\
    \bottomrule
  \end{tabular}
  \caption{Total Effect (TE) of Donald Trump sentiment on ideology label for
    two percentile contrasts. $N = 233$ articles. All five paradigms agree on a
    negative direction; only Fine-tuned GPT-4o-mini's Total Effect reaches 95\%
    significance under either contrast (marked $^*$). The proportion mediated
    through Politics sentiment is reliable for Fine-tuned GPT-4o-mini at
    $\approx 0.60$.}
  \label{tab:te-trump}
\end{table}

All five paradigms agree on the direction of the Trump-sentiment effect: more positive Trump sentiment is associated with a leftward shift in ideology classification. Fine-tuned GPT-4o-mini is the only paradigm whose Total Effect reaches 95\% significance, doing so under both percentile contrasts ($\widehat{\text{TE}} = -0.249$, 95\% CI $[-0.465, -0.033]$ for $Q25 \to Q75$; $\widehat{\text{TE}} = -0.430$, 95\% CI $[-0.802, -0.058]$ for $P10 \to P90$). The other four paradigms produce wide intervals that span zero. Ordering the paradigms by $|\text{TE}|$ on the $Q25 \to Q75$ contrast yields Fine-tuned GPT-4o-mini ($-0.249$) > Human ($-0.199$) > Fine-tuned RoBERTa ($-0.103$) > Llama-3.3-70B ($-0.091$) > GPT-4o-mini baseline ($-0.078$). Holding the architecture constant, fine-tuning therefore amplifies the Total Effect in GPT-4o-mini by roughly $3.2\times$ relative to its zero-shot baseline; this magnitude amplification, rather than a directional disagreement, is the distinguishing causal footprint of fine-tuning at the topic level.

\authortodo{Do the dataset analysis on articles tagged with Trump/Politics -- are they generally mislabelled?}


In every paradigm, the indirect channel through Politics sentiment carries a larger share of the Total Effect than the direct channel: $|\text{NIE}| > |\text{NDE}|$ for all five annotators on both percentile contrasts. For Fine-tuned GPT-4o-mini, the proportion mediated through Politics is reliable at $\approx 0.60$ under both contrasts, and the bootstrap NIE point estimates ($-0.148$ and $-0.262$) are also the largest in absolute value across paradigms. Together with the significant Total Effect, this places the bulk of Fine-tuned GPT-4o-mini's Trump-sentiment signature inside the Politics-mediator pathway, a structural pattern the chapter returns to in Section~\ref{sec:cross-paradigm}.

% ------------------------------------------------------------------------------
\subsection{Multi-Treatment DML Results}
\label{sec:multi-treatment-results}
% ------------------------------------------------------------------------------

The multi-treatment sweep fits a single \texttt{LinearDML} model per dataset with all $k = 11$ qualifying topic sentiments as simultaneous treatments, controlling for topic-presence confounders via PCA-compressed $W$. ATEs are estimated on the full corpus ($N = \num{9830}$), providing substantially higher statistical power than the mediation subsample.

\subsubsection{Continuous Treatment Mode}
\label{sec:continuous-mode}

Table~\ref{tab:multi-continuous} reports the ATE of a one-standard-deviation increase in each topic's sentiment score, holding all other topic sentiments constant.

\begin{table}[ht]
  \centering
  \sisetup{table-format=+1.4}
  \setlength{\tabcolsep}{4pt}
  \small
  \begin{tabular}{l c
      S[table-format=+1.4]
      S[table-format=+1.4]
      S[table-format=+1.4]
      S[table-format=+1.4]
      S[table-format=+1.4]}
    \toprule
    \textbf{Topic} & {\textbf{N}}
      & {\textbf{Human}}
      & {\textbf{FT RoBERTa}}
      & {\textbf{GPT (base.)}}
      & {\textbf{FT GPT}}
      & {\textbf{Llama-3.3}} \\
    \midrule
    Hillary Clinton & 51 & +0.0051 & \cellcolor{blue!25}+0.0220 & +0.0025 & +0.0055 & -0.0053 \\
    Media Bias & 115 & +0.0024 & +0.0203 & -0.0039 & +0.0080 & -0.0092 \\
    FBI & 72 & +0.0182 & -0.0052 & -0.0028 & +0.0008 & -0.0097 \\
    Donald Trump & 233 & -0.0089 & -0.0025 & -0.0007 & -0.0194 & -0.0016 \\
    Media Watch & 58 & -0.0032 & -0.0079 & -0.0063 & -0.0064 & +0.0019 \\
    Russia Probe & 57 & +0.0011 & +0.0045 & -0.0046 & -0.0031 & -0.0041 \\
    Race and Racism & 104 & -0.0013 & +0.0040 & +0.0030 & +0.0061 & +0.0031 \\
    Defense and Security & 76 & +0.0017 & +0.0034 & +0.0038 & -0.0022 & -0.0005 \\
    LGBTQ Issues & 73 & +0.0027 & +0.0021 & -0.0006 & +0.0142 & -0.0009 \\
    Politics & 594 & -0.0035 & -0.0023 & +0.0020 & +0.0102 & -0.0072 \\
    Economy and Jobs & 56 & -0.0001 & +0.0019 & -0.0064 & -0.0064 & -0.0069 \\
    \bottomrule
  \end{tabular}
  \caption{Multi-treatment DML: ATE per topic under continuous encoding (one-SD increase in sentiment score). Cells shaded in blue denote statistically significant results (95\% bootstrap CI excludes zero; $B = \num{2000}$). The full table with 95\% confidence intervals for every estimate is reported in Appendix~\ref{sec:appendix-community-mediation}, Table~\ref{tab:multi-continuous-full}.}
  \label{tab:multi-continuous}
\end{table}

\subsubsection{The Hillary Clinton Signal}
\label{sec:hillary-signal}

Only one cell in the multi-treatment continuous sweep reaches significance: Hillary Clinton sentiment under Fine-tuned RoBERTa, whose ATE narrowly excludes zero ($\widehat{\text{ATE}} = +0.0220$, 95\% CI $[+0.000, +0.035]$). This is the only significant topic-level effect produced by Fine-tuned RoBERTa across any of the topic-level analyses, and it is not recovered by any other paradigm. Outside this isolated finding, no topic is significant under any other annotation paradigm in the multi-treatment continuous analysis.

\authortodo{Do the dataset analysis on articles tagged with Clinton -- are they generally mislabelled?}

% The Hillary Clinton direction under Fine-tuned RoBERTa is counterintuitive: \emph{more positive} Clinton sentiment is associated with \emph{more right-leaning} ideology classification. A plausible explanation is selection bias: the 51 Clinton-tagged articles are disproportionately from 2016--2017, where outlets discuss her in a political-foil context (e.g., framing ``deep state'' criticism via Clinton references). The effect is small and non-significant under every other paradigm (Human $+0.0051$, baseline GPT $+0.0025$, Fine-tuned GPT $+0.0055$, Llama $-0.0053$), and the four LLM-based paradigms cluster within $\pm 0.006$ of zero. The fact that the only significant result comes from a single annotation paradigm strongly suggests the Fine-tuned RoBERTa signal is driven by labelling idiosyncrasies rather than a genuine causal relationship.

% ------------------------------------------------------------------------------
\subsection{Topic-Level Cross-Paradigm Comparison}
\label{sec:cross-paradigm}
% ------------------------------------------------------------------------------

\subsubsection{Agreement Across Paradigms}

Table~\ref{tab:cross-paradigm-summary} consolidates the directional agreement and significance of key causal estimates across all completed paradigm runs.

\begin{table}[ht]
  \centering
  \resizebox{\textwidth}{!}{%
  \begin{tabular}{llcccccccccc}
    \toprule
    \textbf{Analysis} & \textbf{Topic} &
      \multicolumn{2}{c}{\textbf{Human}} &
      \multicolumn{2}{c}{\textbf{FT RoBERTa}} &
      \multicolumn{2}{c}{\textbf{GPT (baseline)}} &
      \multicolumn{2}{c}{\textbf{FT GPT-4o-mini}} &
      \multicolumn{2}{c}{\textbf{Llama-3.3-70B}} \\
    \cmidrule(lr){3-4}\cmidrule(lr){5-6}\cmidrule(lr){7-8}\cmidrule(lr){9-10}\cmidrule(lr){11-12}
    & & Dir. & Sig. & Dir. & Sig. & Dir. & Sig. & Dir. & Sig. & Dir. & Sig. \\
    \midrule
    Mediation (DT trt.)  & Trump $\to$ ideology
      & $-$ & No & $-$ & No & $-$ & No & $-$ & \textbf{Yes} & $-$ & No \\
    Multi-trt.  & Hillary Clinton
      & $+$ & No & $+$ & \textbf{Yes} & $+$ & No & $+$ & No & $-$ & No \\
    Multi-trt.  & Donald Trump
      & $-$ & No & $-$ & No & $-$ & No & $-$ & No & $-$ & No \\
    Multi-trt.  & Politics
      & $-$ & No & $-$ & No & $+$ & No & $+$ & No & $-$ & No \\
    \bottomrule
  \end{tabular}%
  }
  \caption{Summary of cross-paradigm agreement for key ATE estimates.
    ``Dir.'' = direction of ATE (positive/negative).
    ``Sig.'' = statistically significant at $\alpha = 0.05$.
    ``$\sim 0$'' = ATE indistinguishable from zero.}
  \label{tab:cross-paradigm-summary}
\end{table}

\subsubsection{Divergence Analysis}

Two patterns stand out. First, the topic-level Trump-mediation analysis exhibits \textbf{full directional agreement under fine-tuning amplification}: all five paradigms estimate a negative Total Effect on identical articles, but only Fine-tuned GPT-4o-mini's interval excludes zero, and its TE magnitude ($-0.249$) is roughly $3.2\times$ that of the zero-shot GPT-4o-mini baseline ($-0.078$) on the same article subset. Holding architecture constant, fine-tuning therefore alters the magnitude of the causal effect that the labelling model encodes without changing its direction. This is the cleanest topic-level signature attributable to the fine-tuning process.

Second, the multi-treatment continuous analysis produces a single significant cell: Hillary Clinton sentiment under Fine-tuned RoBERTa. No other paradigm recovers this finding, and Fine-tuned RoBERTa recovers no other significant topic. Combined with the fact that the multi-treatment Politics and Trump topics are non-significant under every paradigm, this isolated significance is a strong indicator that the Hillary Clinton signal is an annotation-model artifact rather than a robust causal finding.

% ------------------------------------------------------------------------------
\subsection{Topic-Level Summary}
\label{sec:causal-summary-topic}
% ------------------------------------------------------------------------------

The topic-level analyses produce a narrow but coherent set of significant findings. In the Donald-Trump-as-treatment mediation analysis, all five paradigms agree on a negative direction, but only Fine-tuned GPT-4o-mini reaches significance, with a Total Effect roughly $3.2\times$ the magnitude of the zero-shot GPT-4o-mini baseline on the same article subset. In the multi-treatment continuous analysis, only Hillary Clinton sentiment under Fine-tuned RoBERTa reaches significance, and that finding is not recovered by any other annotation paradigm. Human, GPT-4o-mini (baseline), and Llama-3.3-70B annotations produce no significant topic-level effects across any analysis.

Two distinct empirical patterns are worth separating. The first is a magnitude effect under directional agreement: in the Trump-mediation pathway, fine-tuning a generative LLM amplifies an effect that every other paradigm encodes in the same direction, sharpening it into significance without changing its sign. The second is paradigm-specific isolated significance: in the multi-treatment analysis, no topic produces a statistically significant effect that is consistent across more than one annotation paradigm, and the few significant cells (one per fine-tuned paradigm in the multi-treatment sweep) fail to replicate elsewhere. The first pattern points to fine-tuning as a magnitude amplifier of an underlying shared signal; the second points to paradigm-dependent labelling artefacts at the multi-treatment level.

The primary value of the topic-level analysis lies in demonstrating that identical causal analyses applied to labels from different annotation systems can produce qualitatively different conclusions, both in the magnitude of significant effects (Trump mediation, where fine-tuning amplifies a shared direction) and in the location of multi-treatment signal (where Fine-tuned RoBERTa alone finds Hillary Clinton). As shown in the preceding community-level analysis (Section~\ref{sec:community-level-results}), aggregation across related topics does reveal further paradigm-specific patterns, but these too remain narrow rather than universal.
